% !TEX TS-program = pdflatex
% !TEX encoding = UTF-8 Unicode
%
% Mattia Domizio - CV (EN)
% Layout adapted from Luigi Gonnella's template
%
\documentclass[10pt,a4paper]{article}

\usepackage[a4paper,top=1.4cm,bottom=1.4cm,left=1.6cm,right=1.6cm]{geometry}
\usepackage[T1]{fontenc}
\usepackage[utf8]{inputenc}
\usepackage{lmodern}
\usepackage{microtype}
\usepackage{titlesec}
\usepackage{enumitem}
\usepackage{xcolor}
\usepackage{tabularx}
\usepackage{array}
\usepackage{parskip}
\usepackage{fontawesome5}
\usepackage[hidelinks]{hyperref}

\pagestyle{empty}
\setlength{\parindent}{0pt}

% --- Section formatting ---------------------------------------------------
\titleformat{\section}
  {\Large\bfseries\MakeUppercase}
  {}{0pt}{}[\vspace{-6pt}\hrulefill]
\titlespacing*{\section}{0pt}{10pt}{6pt}

% --- Helpers --------------------------------------------------------------
\newcommand{\entry}[4]{%
  \noindent\begin{tabularx}{\linewidth}{@{}X r@{}}
    \textbf{#1} & #3 \\
    \textit{#2} & \\
  \end{tabularx}\vspace{-4pt}
  \ifx&#4&\else\par#4\fi
}

\newcommand{\simpleentry}[3]{%
  \noindent\begin{tabularx}{\linewidth}{@{}X r@{}}
    \textbf{#1} & #3 \\
    \textit{#2} & \\
  \end{tabularx}\vspace{2pt}
}

\setlist[itemize]{leftmargin=1.2em,topsep=2pt,itemsep=1pt,parsep=0pt}

% --- Header ---------------------------------------------------------------
\begin{document}

\begin{center}
  {\Huge\bfseries Mattia Domizio}\\[2pt]
  {\itshape Computer Engineer \,$|$\, Software Engineer \,$|$\, Data Engineer}\\[6pt]
  \small
  \faEnvelope[regular]\, \href{mailto:mat.dom1609job@gmail.com}{mat.dom1609job@gmail.com}\quad
  \faPhone\, +39 3899014089\quad
  \faLink\, \href{https://koisuji02.github.io/}{koisuji02.github.io}\\[2pt]
  \faLinkedin\, \href{https://linkedin.com/in/mattia-domizio-7a07b628b}{linkedin.com/in/mattia-domizio-7a07b628b}\quad
  \faGithub\, \href{https://github.com/Koisuji02}{github.com/Koisuji02}
\end{center}

% --- Summary --------------------------------------------------------------
\section*{Summary}
Computer Engineering MSc student at Politecnico di Torino with hands-on experience across
\textbf{full-stack web development}, \textbf{distributed systems}, \textbf{data engineering}
and \textbf{low-level system and GPU programming}. Currently working on an \textbf{external thesis at Target Reply}
on a metadata platform. Looking for full-time role as
\textbf{Data/Software Engineer} where I can ship reliable systems and keep learning.

% --- Education ------------------------------------------------------------
\section*{Education}

\entry{MSc in Computer Engineering (LM-32)}{Politecnico di Torino}{2024 -- Present}{%
\begin{itemize}
  \item Exam grades:
    \textbf{Data Science and Database Technology} (30),
    \textbf{Networks Services and Technologies} (29),
    \textbf{Computer Architectures} (30),
    \textbf{Information Systems} (27),
    \textbf{Web Applications I} (30L),
    \textbf{Software Engineering} (30),
    \textbf{Formal Languages and Compilers} (28),
    \textbf{System and Device Programming} (28),
    \textbf{Distributed Systems Programming} (29),
    \textbf{Software Engineering II} (30L),
    \textbf{Information Systems Security} (27),
    \textbf{GPU Programming} (29).
\end{itemize}}

\entry{BSc in Computer Engineering (L-08)}{Politecnico di Torino}{2021 -- 2024}{%
\textbf{107/110}}

\entry{Scientific High School Diploma}{Liceo Scientifico Cattaneo, Turin}{2021}{%
\textbf{100/100}}

\simpleentry{Cambridge B2 First Certificate}{English language certification}{2020}

% --- Work Experience ------------------------------------------------------
\section*{Work Experience}

\entry{Data Engineer Intern -- External Thesis Project}{Target Reply S.r.l.}{02/2026 -- Present}{%
\begin{itemize}
  \item Working on the master's thesis project taking part in the development of a metadata governance platform.
\end{itemize}}

\entry{University Tutor}{Politecnico di Torino}{02/2025 -- 07/2025}{%
\begin{itemize}
  \item Supported undergraduate students with lab exercises and code reviews on
        the OOP course (Java).
\end{itemize}}

\entry{Weekend Fruit \& Vegetable Clerk}{Gabrielis S.r.l.}{11/2023 -- 02/2024}{}

\entry{Election Poll Worker}{Comune di Torino}{09/2022}{}

\entry{Private Tutor}{Academic support and homework assistance}{2020 -- 2024}{}

\entry{Assistant Coach}{A.S.D. Pozzomaina}{10/2017 -- 06/2019}{}

% --- Skills ---------------------------------------------------------------
\section*{Skills}

\textbf{Languages}\\
Python, Rust, C, C++, Java, TypeScript, JavaScript, Bash, Powershell, GDScript, CUDA, SQL, HTML5, CSS3, Assembly Script (ARM, x86-64, RISC-V, 32), Scala, LaTeX

\vspace{4pt}
\textbf{Frontend \& UI}\\
React, React Router, Vite, Next.js, Bootstrap, TailwindCSS, Qt

\vspace{4pt}
\textbf{Backend \& APIs}\\
Node.js, Express, Flask, Django, OpenAPI / Swagger, REST, JWT, Postman, .NET, Spring, GraphQL, gRPC, HTTP/2, WebSockets, MQTT, Kafka

\vspace{4pt}
\textbf{Data \& Databases}\\
PostgreSQL, MySQL, SQLite, MongoDB, Redis, Apache Hadoop, Apache Spark, DataHub,
Power BI, Pandas, NumPy, Matplotlib, PyTorch, TensorFlow, Databricks, Snowflake, dbt, Elasticsearch, Qlik, MicroStrategy

\vspace{4pt}
\textbf{Cloud \& DevOps}\\
Docker, Kubernetes, Git, GitHub, GitLab, GitHub Actions, AWS, Azure, CMake, Helm, NPM, SonarQube

\vspace{4pt}
\textbf{Game Dev \& Design}\\
Godot Engine, Figma, GIMP, NVIDIA tooling

\vspace{4pt}
\textbf{OS \& Tooling}\\
Linux, Windows, Arch, Nix, WireGuard

% --- Additional -----------------------------------------------------------
\section*{Additional}

\textbf{Languages}
\begin{itemize}
  \item Italian (Native)
  \item English (\textbf{B2 First Cambridge})
\end{itemize}

% --- Projects -------------------------------------------------------------
\section*{Projects}

\entry{Godot Shader Pipeline Warm-Up}{Godot Engine, GDScript, GDShader, Graphics Pipeline, Benchmarking}{2025 -- 2026}{%
\begin{itemize}
  \item Designed and implemented a \textbf{deterministic shader pipeline warm-up workflow}
        in Godot to reduce first-use stutter at runtime.
  \item Ran benchmarks comparing cold vs.\ warmed-up runs and documented results in a
        dedicated \textbf{report.pdf}.
  \item \textbf{GitHub}: \href{https://github.com/Koisuji02/GodotShaderWarmup}{link}
\end{itemize}}

\entry{Participium -- Civic Reporting Platform}{Full-stack Web (TypeScript, Python, React, JavaScript, HTML, CSS), Software Engineering (Swagger / OpenAPI, Docker, Bash, Redis), Agile (YouTrack)}{2026}{%
\begin{itemize}
  \item Civic reporting platform for the city of Turin: citizens submit reports on urban
        issues and municipal officers triage and manage them.
  \item Exam project for \textbf{Software Engineering II} at Politecnico di Torino.
  \item \textbf{GitHub}: \href{https://github.com/Koisuji02/Participium}{link}
\end{itemize}}

\entry{GeoControl}{REST APIs, TypeScript, React, JavaScript, HTML, CSS, Docker, Swagger / OpenAPI}{2025}{%
\begin{itemize}
  \item Monitoring system for heterogeneous sensors (position, temperature, pressure,
        humidity, gas concentration and other geo-environmental variables).
  \item Designed the REST API surface and the underlying data model; full documentation
        shipped under the \texttt{/doc} folder.
  \item \textbf{GitHub}: \href{https://github.com/Koisuji02/GeoControl}{link}
\end{itemize}}

\entry{Ruggine -- Rust Chat App}{Rust, Tokio, WebSockets, React, JavaScript, HTML, CSS}{2025}{%
\begin{itemize}
  \item Built a \textbf{secure messaging system} in Rust with private and group chats.
  \item Exam project for \textbf{System and Device Programming} (\textit{Programmazione di Sistema} course) at Politecnico di Torino.
  \item \textbf{GitHub}: \href{https://github.com/Koisuji02/RuggineChatRust}{link}
\end{itemize}}

\entry{Stuff Happens}{React, Node.js, Express, SQLite, JWT, JavaScript, HTML, CSS}{2025}{%
\begin{itemize}
  \item Web-based card game where players guess the position of a new ``bad luck''
        card relative to three known cards, scoring on a custom \textit{bad luck index}.
  \item Features \textbf{user authentication}, game history tracking and a
        travel-mishap themed deck (``Unlucky Holidays'').
  \item \textbf{GitHub}: \href{https://github.com/Koisuji02/StuffHappens}{link}
\end{itemize}}

\entry{Film Manager Service}{REST APIs, Swagger / OpenAPI, JSON Schema, JavaScript, Postman}{2024 -- 2025}{%
\begin{itemize}
  \item REST service to manage a film catalogue, with an \textbf{OpenAPI/Swagger}
        specification, \textbf{JSON Schemas} for payload validation and a Postman
        collection covering the main flows.
  \item Exam Call 1 project (\textbf{Distributed Systems}).
  \item \textbf{GitHub}: \href{https://github.com/Koisuji02/FilmManagerService}{link}
\end{itemize}}

\entry{PAC-MAN on LPC1768}{ARM Assembly, C, Embedded C}{2023 -- 2024}{%
\begin{itemize}
  \item Ported PAC-MAN to the \textbf{LPC1768} board using ARM \textbf{assembly} and C,
        targeting the on-board peripherals.
  \item Exam project for \textbf{Computer Architectures} at Politecnico di Torino.
  \item \textbf{GitHub}: \href{https://github.com/Koisuji02/PAC-MAN}{link}
\end{itemize}}

\end{document}
